\chapter{Profil Perusahaan dan Kegiatan MBKM}

\section{Profil Perusahaan}
PT Vortex Buana Edumedia merupakan perusahaan teknologi dan transformasi digital yang berpusat di Yogyakarta. Vortex berfokus pada pengembangan solusi perangkat lunak end-to-end untuk kebutuhan pendidikan dan bisnis, baik melalui produk internal maupun layanan pengembangan digital. Dalam ekosistem produknya, Vortex mengembangkan Neon, yang mencakup Neon Uji dan Neon Belajar. Neon Uji menyediakan kuis on-demand, tryout, laporan hasil yang rinci, serta fitur pemantauan progres peserta yang mendukung proses evaluasi pembelajaran secara lebih terstruktur. Selain pengembangan produk, Vortex juga menyediakan layanan pengembangan web dan media digital yang mendukung kebutuhan transformasi digital mitra dan pengguna.

\section{Penjelasan Program Magang}
Selama program MBKM Independen, saya melaksanakan magang selama tiga bulan, terhitung sejak 3 November 2025 hingga 3 Februari 2026, sebagai Frontend Developer pada divisi Product di PT Vortex Buana Edumedia. Pada kegiatan ini, saya terlibat dalam pengembangan beberapa produk internal yang berfokus pada penguatan ekosistem Neon serta penyempurnaan pengalaman pengguna pada sistem yang sedang berjalan.

Proyek pertama yang saya kerjakan adalah Neutron Office, yaitu sistem integrasi yang menghubungkan rekam belajar dengan nilai dari Neon Uji. Sistem ini dirancang untuk menyatukan data pembelajaran dan penilaian agar dapat dikelola secara lebih rapi dan mudah diakses. Kontribusi saya pada proyek ini berfokus pada pengembangan antarmuka frontend, penyesuaian alur interaksi pengguna, serta implementasi komponen yang mendukung integrasi data dengan tampilan yang jelas dan mudah digunakan.

Proyek kedua adalah Neon Uji SaaS, yaitu platform asesmen yang digunakan untuk kebutuhan kuis, tryout, dan evaluasi pembelajaran. Pada proyek ini, saya berperan dalam perbaikan bug, peningkatan kualitas antarmuka, serta penyempurnaan pengalaman pengguna agar sistem lebih stabil, konsisten, dan nyaman digunakan pada berbagai skenario penggunaan. Pengerjaan ini menuntut ketelitian tinggi karena perubahan kecil pada antarmuka dapat memengaruhi alur penggunaan secara keseluruhan.

Proyek ketiga adalah Neon Interfaces, yaitu dokumentasi design library yang berfungsi sebagai referensi komponen antarmuka bagi tim desain dan pengembangan. Pada proyek ini saya membantu menyusun dokumentasi komponen, membangun basis design library yang mengacu pada shadcn/ui, serta menjaga konsistensi pola UI agar dapat digunakan kembali pada proyek-proyek internal lain. Pengembangan ini mendukung standarisasi tampilan sekaligus mempercepat implementasi antarmuka di tim Product.

Selama magang, saya menggunakan Nuxt 3, Vue 3, Storybook, MongoDB, Tailwind CSS, dan perangkat pendukung lainnya. Pengalaman tersebut memperkuat kemampuan saya dalam membangun antarmuka modular, bekerja dengan design system, melakukan penyesuaian komponen secara konsisten, serta berkolaborasi dengan tim lintas fungsi dalam lingkungan pengembangan produk.

Kegiatan magang ini berkaitan erat dengan beberapa mata kuliah MBKM, khususnya Internship: Pengembangan Fitur dan Modul Proyek (MII21-3011), Internship: Pengujian Unit dan Modul Proyek (MII21-3015), serta Soft Skill: Kemampuan Berkomunikasi (MII21-4010). Keterkaitan tersebut terlihat dari proses penerjemahan kebutuhan produk menjadi fitur frontend, validasi perubahan melalui pengujian, dan koordinasi aktif dengan tim Product, desain, serta pengembang lain dalam menyelesaikan pekerjaan.

Pengalaman magang ini memberikan pemahaman yang lebih nyata mengenai proses pengembangan perangkat lunak di lingkungan industri, terutama pada aspek konsistensi antarmuka, pengelolaan komponen yang dapat digunakan ulang, dan kolaborasi tim. Karena sebagian hasil pekerjaan berada dalam sistem internal perusahaan, dokumentasi implementasi dan hasil akhir dibahas lebih lanjut pada bab-bab berikutnya serta didukung oleh tangkapan layar pada bagian lampiran.